% ECONOMICS_PROBLEM: {"id": "H77-307", "title": "When to decentralize government", "jel_code": "H77", "source_id": "OP-0386"}
% !TeX program = xelatex
\documentclass[11pt,a4paper]{article}
\usepackage[margin=23mm]{geometry}
\usepackage{fontspec}
\setmainfont{Latin Modern Roman}
\usepackage{amsmath,amssymb,mathtools}
\usepackage{xcolor,enumitem,booktabs,longtable,array,xurl,hyperref}
\definecolor{muted}{HTML}{53636C}
\hypersetup{colorlinks=true,urlcolor=blue,linkcolor=blue}
\setlength{\parindent}{0pt}
\setlength{\parskip}{5pt}
\newcommand{\R}{\mathbb R}
\newcommand{\E}{\mathbb E}
\newcommand{\N}{\mathbb N}
\newcommand{\ind}{\mathbf 1}
\newcommand{\Var}{\operatorname{Var}}
\newcommand{\Cov}{\operatorname{Cov}}
\newcommand{\supp}{\operatorname{supp}}
\DeclareMathOperator*{\argmin}{arg\,min}
\DeclareMathOperator*{\argmax}{arg\,max}
\newcommand{\entrymeta}[3]{{\sffamily\small\color{muted}\textbf{\texttt{#1}}\quad|\quad #2\par #3\par}}
\newcommand{\Problemref}[1]{\texttt{#1}}
\newcommand{\JELref}[2]{\texttt{#2} (JEL \texttt{#1})}
\begin{document}
\section*{When to decentralize government}
\textbf{Problem ID:} \texttt{H77-307}\quad \textbf{JEL:} \texttt{H77}\par
% BEGIN ECONOMICS STATEMENT
\label{problem:OP-0386}
{\sffamily\small\color{muted}Original subject group: Political economy and institutions\par}
\label{definitions:problem:30007668}
\textbf{Audit classification:} Published research agenda. Full review checked. Adaptation/coordination tradeoffs are explicit research work; authority vectors, constraints and causal rules are editorial.
\subsection*{Definitions and precise research target}
\textbf{Complete editorial problem.} Consider regions \(j\) within a specified state providing a public service. Governance rule \(d\in[0,1]\) specifies the degree of local authority over budgets, personnel, and service design, separately from the source of financing. Let \(Q_j(d)\) be service quality, \(K_j\) prereform local administrative capacity, \(I_j\) local information, and \(E_j\) prereform elite influence. Let \(C(d)\) be total cost and \(V(d)=\sum_j\omega_jQ_j(d)\) a declared population-weighted objective. Estimate heterogeneous policy effects \(E[Q_j(d)-Q_j(0)\mid K_j,I_j,E_j]\) and characterize the feasible rule maximizing \(V\) subject to a budget and explicit inequality constraints. Decentralization can improve local adaptation while reducing coordination, so include cross-region spillovers and national distributional effects. Information technology may change the value of local information and must be incorporated into the policy environment. Use credible reform or assignment variation and avoid treating selected high-capacity localities as representative. The cited review expressly identifies this organizational tradeoff as a research area. The proposed causal contrasts and constrained authority rule make the broad question precise without assuming decentralization has the same effect for every service or state.

\textbf{Definitions and admissible scope.} Authority is naturally a vector \(d=(d_b,d_h,d_s)\in[0,1]^3\) over budgets, hiring, and service design; a scalar \(d\) is valid only after fixing a one-dimensional path through that vector. Financing and equalization transfers are separate policy variables. Regions have weights \(\omega_j\ge0\), summing to one, and outcomes \(Q_j(\mathbf d)\) depend on the full regional assignment because of spillovers. Set \(V(\mathbf d)=E[\sum_j\omega_jQ_j(\mathbf d)]\). A regional rule maps prereform \((K_j,I_j,E_j)\) into authority; admissibility requires measurable mapping, budget, and specified access or inequality constraints. Local information and elite influence indices have declared measurement rules. Estimating a region's effect conditional on baseline capacity does not identify a uniform national optimum. In the scalar base case fix a known path \(h:[0,1]\to[0,1]^3\) with \(h(0)\) the central-authority benchmark, and read \(Q_j(d)\) as outcomes under the full assignment of \(h(d)\). In the regional-rule extension let \(\mathbf d=(d_1,\ldots,d_J)\), each \(d_j\in[0,1]^3\). A feasible rule is a measurable map of baseline region characteristics into \(\mathbf d\), with fixed financing/equalization rules, expected cost at most \(B\), and explicit constraints such as \(E[Q_j]\ge\underline Q_j\) for every region. Finite-mean quality, citizen weights, assignment/spillover laws, and response-model restrictions are inputs. Identify the feasible value set rather than interpreting selected-locality estimates as the national rule's outcome.
\subsection*{Variants and assumption changes}
\begin{enumerate}
\item Administrative versus fiscal authority (editorial): vary personnel delegation while national finance and equalization remain fixed, then vary finance separately. Estimate whether the complementarity differs by baseline capacity.
\item Information technology and coordination (source-motivated extension): compare central real-time information systems with local discretion when service spillovers are large. Measure total system quality and cross-region inequality, not only treated-region gains.
\end{enumerate}
\subsection*{References and source audit}
\textbf{Support and access.} Full review checked. Adaptation/coordination tradeoffs are explicit research work; authority vectors, constraints and causal rules are editorial. \textbf{Locator:} Section 4.1, p. 413, paragraph on hierarchy and decentralization.
\begin{enumerate}\raggedright
\item\hypertarget{definitions:source:30007668:1}{} Timothy Besley; Robin Burgess; Adnan Khan; Guo Xu. 2022. \emph{Bureaucracy and Development}. Section 4.1, printed p. 413, paragraph beginning "Although a hierarchical organization".
\url{https://www.robinburgess.com/s/Bureaucracy_and_Development.pdf}.
DOI: \url{https://doi.org/10.1146/annurev-economics-080521-011950}.
\end{enumerate}

\subsection*{Common conventions: Source mathematical conventions}

These conventions apply to each entry unless explicitly replaced.

\textbf{Models and identification.} A primitive model \(m\) specifies all probability laws, feasible choices, information, equilibrium and measurement rules used by its target. The admissible class \(\mathcal M\) is fixed independently of outcomes. If \(P_O(m)\) is its observable probability law and \(\tau(m)\) the target, its identified set at observable law \(P\) is
\[
 \mathcal I_\tau(P)=\{\tau(m):m\in\mathcal M,\ P_O(m)=P\}.
\]
Point identification means that this set is a singleton. An empty set signals incompatibility of data and assumptions. Coordinate bounds need not describe the joint set. Bounds are sharp only if every retained target is attainable or arbitrarily approachable by compatible models. Finite bounds require explicit restrictions on unobserved quantities; unrestricted confounding, hidden wealth, extrapolation or arbitrary equilibrium selection can yield uninformative sets.

\textbf{Causal interventions.} A potential outcome \(Y_i(p)\) is the outcome under a complete policy regime \(p\), including timing, financing, information and market spillovers. The target population and units are fixed before treatment. With interference, use the complete assignment vector or a justified exposure mapping; individual no-interference notation alone cannot identify an economy-wide intervention. A difference of mean potential outcomes does not require the cross-world joint distribution. Conditioning on post-treatment migration, survival, enrollment or employment changes the estimand and needs separate justification.

\textbf{Statistical statements.} An estimator, confidence set, forecast or test specifies a sampling law, model class and loss/coverage criterion. Uniform \((1-\alpha)\)-coverage means \(\inf_{m\in\mathcal M}\Pr_m\{\tau(m)\in C_n\}\geq1-\alpha\), or an explicitly asymptotic version. The whole parameter space trivially has coverage, so informativeness requires a stated width, risk or power objective. A forecasting improvement is relative to a named benchmark, information set and evaluation population.

\textbf{Equilibrium and optimization.} An equilibrium comprises feasible plans, optimality, beliefs where needed, budget constraints and all market/resource clearing. The primitives or a stated benchmark define each correspondence \(\mathcal E(p;m)\). Nonemptiness is assumed or proved, never inferred just from compact policy sets. A selected-equilibrium target uses a declared selection rule; a robust target quantifies over all permitted equilibria. Use suprema or infima unless attainment is established. An \(\varepsilon\)-optimal policy is feasible and within \(\varepsilon\) of the finite optimum value. Report an empty feasible set separately.

\textbf{Welfare and accounting.} Normative utility, interpersonal normalization, social weights, numeraire, discounting and the use of public receipts are primitives. Behavioral utility need not coincide with the welfare criterion. Household welfare includes ownership income and rents once; adding profits again to compensating variation generally double-counts them. A monetary equivalent is defined by an explicit transfer and price convention, with monotonicity and range conditions. Average effects, mechanism contributions, transfers and welfare losses are different targets. Infinite sums require convergence and dynamic feasible plans require terminal or no-Ponzi conditions.

\textbf{Source versions.} A working-paper date, revision date and journal publication year are distinguished. A DOI for a journal version is not automatically the DOI of an earlier manuscript. Locators refer to the stated version and printed pages where specified. A metadata-only check does not verify a claim about a full-text conclusion. Every entry's source subsection narrows the provenance claim accordingly.
% END ECONOMICS STATEMENT
\end{document}
